\documentclass[10pt,a4paper]{amsart}
\usepackage[margin=19mm]{geometry}
\usepackage{amsmath,amssymb,mathtools}
\usepackage[scr=rsfs,cal=euler]{mathalfa}
\usepackage{xfrac}
\usepackage[unicode,hidelinks,bookmarks=false]{hyperref}
\newcommand{\ie}{i.e.\ }
\newcommand{\cf}{cf.\ }
\newcommand{\resp}{respectively\ }
\newcommand{\Subsection}{\subsection}
\providecommand{\nobreakdash}{\nobreak\hbox{-}}
\setlength{\emergencystretch}{2em}
\multlinegap0pt
\usepackage{amssymb}
\usepackage{mathtools}
\newtheorem{lemma}{Lemma}
\newcommand{\B}{\mathcal{B}}
\newcommand{\na}{{\nabla}}
\newcommand{\vol}{{\rm{vol}}}
\newcommand{\vphi}{\varphi}
\title{Quantitative Stability for Minimizing Yamabe Metrics with minimal boundary}
\author{Runze Lin, Bao Yu}
\date{Source: arXiv:2605.17191v2 (CC BY 4.0)}
\begin{document}
\maketitle

\noindent\emph{Source and licence.} Quantitative Stability for Minimizing Yamabe Metrics with minimal boundary. Version arXiv:2605.17191v2 (CC BY 4.0), distributed by arXiv under the Creative Commons Attribution 4.0 International licence, http://creativecommons.org/licenses/by/4.0/, as recorded in the arXiv licence metadata for this exact version and shown on the abstract page https://arxiv.org/abs/2605.17191v2. Full licence notice: \texttt{licenses/arxiv-2605.17191-CC-BY-4.0.txt}.\par\smallskip

\noindent\textit{Editorial scope.} Counted unit: the article's lemma lem:bm (source0.tex lines 253--279). The article's complete original proof of it is delivered with it (source0.tex lines 280--345, one \texttt{proof} environment of 66 lines). No mathematics is written, repaired or re-derived: the statement and the proof are the article's own lines, in the article's own notation, and no retained line is substituted or re-flowed.  No further source-local context is needed: every cross-reference of every retained line resolves inside the two delivered spans.  Nothing was omitted from any retained span, so the package carries no omission record.  No retained line contains a citation, so the wrapper carries no bibliography and no bibliography entry.  No retained line contains a figure environment, an image-inclusion command or drawing code, so no image is delivered and the package declares no asset dependency.  Editorial formatting only, in addition: an AMS \texttt{amsart} preamble that restates the article's own declarations and macros used by the retained lines, namely the environments \texttt{lemma} on separate counters, together with the standard packages those lines need.  The retained environments print in this document as Lemma 1 in order of appearance, whereas the article prints the same items with the numbering of its own full document.  The complete native source of the article as distributed in the arXiv e-print for this exact version is bundled unchanged as \texttt{sources/200780/source0.tex}.
\begin{lemma}\label{lem:bm}
    For every $u \in \B$, the first and second variations of $Q$ on $\B$ are given by
    \begin{equation}\label{e:firstvariation}
    \nabla_{\B} Q(u)[\varphi] = 2\int_{M} \left(-\Delta u + c_nR_gu\right)\pi_{T_u\B}\varphi\, d\vol_g\, +2\int_{\partial M} \left( \frac{\partial u}{\partial \nu}+\frac{n-2}{2}h_g u   \right)\pi_{T_u\B}\varphi d\sigma_g ,
    \end{equation}
	\begin{equation}\label{e:secondvariation1}
	\begin{aligned}
	&\frac{1}{2}\nabla_\B^2 Q(u)[\varphi,\eta]
	 = \int_M \left\{  \,\na \pi_{T_u\B}\varphi\,\cdot \na \pi_{T_u\B}\eta + c_n R_g\,(\pi_{T_u\B}\varphi) \,(\pi_{T_u\B}\eta) \right\}\, d\vol_g\\ &+\frac{n-2}{2}\int_{\partial M}h_g(\pi_{T_u\B}\phi) (\pi_{T_u\B}\psi) d\sigma_g-(2^*-1)Q(u)\int_{M} u^{2^*-2}\, (\pi_{T_u\B}\varphi) \,(\pi_{T_u\B}\eta)\, d\vol_g,\end{aligned}
	\end{equation}
 for all $ \varphi,\eta\in W^{1,2}(M)$.	We will often omit the projection maps when we are doing computations with $\nabla_{\B}^2Q$.

 Moreover, the following two continuous properties of the Hessian operator hold.

 Let $\mathcal{L}_v \phi = L_g \phi -(2^*-1)Q(v)v^{2^*-2} \phi $, for $u,w \in C^{2,\alpha} \cap \B$ sufficiently close, there exists a modulus of continuity $\omega$, with $\omega(r) \to 0$ as $r\to 0$, such that  for all $\eta \in C^{2,\alpha}(M)$,
 \begin{equation}\label{prop1}
     ||(\mathcal{L}_u \pi_{T_u{\B}} \eta, B_u \pi_{T_u{\B}} \eta )-(\mathcal{L}_w \pi_{T_w{\B}} \eta, B_w \pi_{T_w{\B}} \eta )||_{C^{0,\alpha}(M) \times C^{1,\alpha}(\partial M)} \leq \omega(||u-w||_{C^{2,\alpha}(M)})||\eta||_{C^{2,\alpha}(M)}.
 \end{equation}

 For $u,w \in \B$ sufficiently close in $W^{1,2}$, there exists a modulus of continuity $\omega$, with $\omega(r) \to 0$ as $r\to 0$, such that  for all $\vphi, \eta \in W^{1,2}(M)$,
 \begin{equation}\label{prop2}
     |\na ^2 Q_\B(u)[\vphi,\eta]-\na ^2 Q_\B(w)[\vphi,\eta]|\leq w(||u-w||_{W^{1,2}(M)})||\vphi||_{W^{1,2}(M)}||\eta||_{W^{1,2}(M)}.
 \end{equation}



\end{lemma}

\begin{proof} 
	  Let us denote
	$$
	\mathcal E(u):=\int_M  |\nabla u|^2+ c_n R_gu^2\,d\vol_g\,+\frac{n-2}{2}\int_{\partial M} h_g u^2 d\sigma_g,
	$$ 
	and observe that if $u\in \B$, then $\mathcal E(u)=Q(u)$. For $u\in \B$ and $\varphi \in W^{1,2}(M)$,  we can compute the first variation of $Q$ at points of $\B$:
	\begin{align}
	\nabla Q(u)[\phi] &:= \frac{d}{dt}(Q (u + t\phi))\Big|_{t=0}\\ 
	&=\bigg( [\vol(M, g_{u+t\phi})]^{-2/2^*}\, 2\big[\int_M\left(( \nabla u\cdot\nabla \phi+\frac{t}{2}\,|\nabla \phi|^2)+c_n R_g(u\,\phi+\frac{t}{2}\phi^2)\right)\,d\vol_g \notag\\
    &+ \frac{n-2}{2}\int_{\partial M}h_g(u\phi+\frac{t}{2}\phi^2)d\sigma_g \big] \notag \\
	 &  -2\,\left[\vol(M, g_{u+t\phi}) \right]^{-(2+2^*)/2^*}\, \left[\int_M (u+t\phi)^{2^*-1} \,\phi \,d\vol_g\right]  \,\mathcal E(u+t\phi)\bigg)\Big|_{t=0} \notag 
	\\
	&= 2 \int_{M} \left(-\Delta u + c_nR_gu-Q(u)u^{2^*-1}\right)\phi\, d\vol_g\, + 2\int_{\partial M} \left( \frac{\partial u}{\partial \nu}+\frac{n-2}{2}h_g u   \right)\phi d\sigma_g.
    \end{align} 
	In particular, when restricted to the tangent space of $\B$, we have
	$$
	\nabla_{\B} Q(u)[\varphi] = 2\int_{M} \left(-\Delta u + c_nR_gu\right)\pi_{T_u\B}\varphi\, d\vol_g\, +2\int_{\partial M} \left( \frac{\partial u}{\partial \nu}+\frac{n-2}{2}h_g u   \right)\pi_{T_u\B}\varphi d\sigma_g .
	$$

	Differentiating \eqref{e:firstvariation}, we obtain
   \begin{align}
       &\nabla^2Q(u)[\phi,\psi]\\
       =& [\vol(M, g_{u})]^{-2/2^*}2\left[\int_M \left(\na \phi \cdot \na \psi +c_n R_g\phi \psi  \right)d\vol_g +\frac{n-2}{2}\int_{\partial M}h_g\phi \psi d\sigma_g    \right] \notag \\
       &-2[\vol(M, g_{u})]^{-2/2^*-1}\left[\int_M u^{2^*-1}\psi d\vol_g  \right]2\cdot\left(\int_M(\na u \cdot \na \phi +c_nR_g u \phi +\frac{n-2}{2}\int_{\partial M}h_g u \phi d\sigma_g \right) \notag \\
       &+2\cdot(2+2^*)[\vol(M, g_{u})]^{-2/2^*-2}\int_M u^{2^*-1}\psi d\vol_g\int_M u^{2^*-1}\phi d\vol_g \mathcal E(u)\notag \\
       &-2[\vol(M, g_{u})]^{-2/2^*-1}(2^*-1)\int_M u^{2^*-2} \phi \psi d\vol_g \mathcal E(u) \notag \\
       &-2[\vol(M, g_{u})]^{-2/2^*-1}\int_M u^{2^*-1} \phi d\vol_g \frac{d}{dt}\large|_{t=0} \mathcal E(u+t\psi) \notag \\
       =&2\left[\int_M \left(\na \phi \cdot \na \psi +c_n R_g\phi \psi  \right)d\vol_g +\frac{n-2}{2}\int_{\partial M}h_g\phi \psi d\sigma_g    \right] -2(2^*-1)Q(u)\int_Mu^{2^*-2}\phi \psi d\vol_g \notag \\
       &-2\int_M u^{2^*-1} \phi d\vol_g \frac{d}{dt}\large|_{t=0} \mathcal E(u+t\psi)-2\int_M u^{2^*-1} \psi d\vol_g \frac{d}{dt}\large|_{t=0} \mathcal E(u+t\phi)\notag \\
       &+2\cdot (2+2^*) Q(u) \int_M u^{2^*-1}\psi d\vol_g\int_M u^{2^*-1}\phi d\vol_g. \notag 
   \end{align}
   In particular, when restricted to the tangent space of $\B$, we have
   \begin{align}
       &\nabla_\B^2Q(u)[\phi,\psi]=2\bigg[\int_M \left(\na \pi_{T_u\B}\phi \cdot \na \pi_{T_u\B}\psi +c_n R_g\pi_{T_u\B}\phi \pi_{T_u\B}\psi  \right)d\vol_g\\
       &+\frac{n-2}{2}\int_{\partial M}h_g\pi_{T_u\B}\phi \pi_{T_u\B}\psi d\sigma_g\bigg]
       -2(2^*-1)Q(u)\int_Mu^{2^*-2}\pi_{T_u\B}\phi \pi_{T_u\B}\psi d\vol_g. \notag
   \end{align}
   To conclude the proof, let us check the continuity. From the definition of $\pi_{T_u{\B}}$, let $a_u(\eta)= \int_M u^{2^* -1} \eta d \vol_g$, we have 
   \begin{align*}
       \pi_{T_u{\B}} \eta -\pi_{T_w{\B}} \eta &= -a_u(\eta)u +a_w(\eta) w\\
       &=-a_u(\eta) (u-w) -(a_u(\eta) -a_w(\eta))w.
   \end{align*}
   It is easy to show $||\pi_{T_u{\B}} \eta -\pi_{T_w{\B}} \eta||_{C^{2,\alpha}(M)} \leq \omega(||u-w||_{C^{2,\alpha}(M)})||\eta||_{C^{2,\alpha}(M)}.$
   Now compute 
   \begin{align} \label{Prop1}
       &\mathcal{L}_u \pi_{T_u{\B}} \eta -\mathcal{L}_w \pi_{T_w{\B}} \eta 
       =L_g (\pi_{T_u{\B}} \eta -\pi_{T_w{\B}} \eta) -(2^*-1)[Q(u)u^{2^*-2}\pi_{T_u{\B}} \eta -Q(w)w^{2^*-2}\pi_{T_w{\B}} \eta] \notag\\
       =&L_g (\pi_{T_u{\B}} \eta -\pi_{T_w{\B}} \eta) -(2^*-1)(Q(u)u^{2^*-2}-Q(w)w^{2^*-2})\pi_{T_u{\B}} \eta - (2^*-1)Q(w)w^{2^*-2}(\pi_{T_u{\B}} \eta -\pi_{T_w{\B}} \eta).
   \end{align}
   Combined with the continuity of $L_g: C^{2,\alpha}(M) \to C^{0,\alpha}(M)$ and $B_g : C^{2,\alpha}(M) \to C^{1,\alpha}(\partial M) $, we can verify \eqref{prop1}.

   For the $W^{1,2} $ continuity, by the Sobolev embedding $W^{1,2} (M) \hookrightarrow L^{2^*}(M)$,
   we get 
   \begin{equation*}
       |a_u(\eta)| \leq ||u||_{L^{2^*}}^{2^*-1} ||\eta ||_{L^{2^*}} \leq C ||\eta ||_{W^{1,2}(M)}.
   \end{equation*}
   Moreover, the map $u \mapsto u^{2^*-1}$ is continuous from $L^{2^*}$ to $L^{\frac{2^*}{2^*-1}
   }$ on nonnegative functions, hence
   \begin{equation*}
       |a_u(\eta)-a_w(\eta)|\leq ||u^{2^*-1}-w^{2^*-1}||_{L^{\frac{2^*}{2^*-1}}}||\eta||_{L^{2^*}}
       \leq \omega(||u-w||_{W^{1,2}(M)}) ||\eta ||_{W^{1,2}(M)}.
   \end{equation*}
   It follows that $||\pi_{T_u{\B}} \eta -\pi_{T_w{\B}} \eta||_{W^{1,2}(M)} \leq \omega(||u-w||_{W^{1,2}(M)})||\eta||_{W^{1,2}(M)}.$
   Then, the same argument as \eqref{Prop1}, we can verify \eqref{prop2}.
   
\end{proof}

\end{document}
